%==================================================================================
\section{Load Matrix with .avs}
The function \texttt{load\_avs} reads an AVS file and returns a matrix of pixels in memory.

As explained earlier, the AVS format stores the image width and height in the first eight bytes of the file, followed by the pixel data. The function reads this information and allocates memory for a matrix of pixels.

Given that we are working in a little-endian architecture, the function must convert the big-endian bytes of the header to little-endian integers.

First of all, we have to define a structure to represent a pixel in memory:

\begin{lstlisting}
typedef struct {
    unsigned char a; // alpha channel
    unsigned char r; // red channel
    unsigned char g; // green channel
    unsigned char b; // blue channel
} Pixel;
\end{lstlisting}

This structure allows us to store the four channels of a pixel in a single object. 

To start loading the image, we open the AVS file in binary read mode, then we read the first four bytes to obtain the image width, and the next four bytes to obtain the image height. We convert these bytes from big-endian to little-endian integers.

\begin{lstlisting}
    FILE *file = fopen(filename, "rb"); // read file as binaries

    if (file == NULL) {
        printf("No file found.\n");
        return NULL;
    }

    // define a 4 bytes buffer to read the header
    unsigned char buffer[4];

    // read first 4 bytes of the header -> width
    fread(buffer, 1, 4, file);

    // arrange big endian bytes to little endian
    width = ((int)buffer[0] << 24) |
            ((int)buffer[1] << 16) |
            ((int)buffer[2] << 8) |
            buffer[3];

    // read the follwing 4 bytes -> height
    fread(buffer, 1, 4, file);

    // arrange header big endian bytes to little endian
    height = ((int)buffer[0] << 24) |
            ((int)buffer[1] << 16) |
            ((int)buffer[2] << 8) |
            buffer[3];
\end{lstlisting}

After obtaining the dimensions, we allocate memory for a matrix of pixels with the specified width and height. 

\begin{lstlisting}
    Pixel **image = malloc(height *sizeof(Pixel *));

    for (int y=0; y < height; y++) {
        // for every y in height, create row
        image[y] = malloc(width * sizeof(Pixel));
    }
\end{lstlisting}

We then read the pixel data from the file, storing each pixel in the corresponding position in the matrix.

\begin{lstlisting}
    for (int y = 0; y < height; y++) {
        for (int x = 0; x < width; x++) {
            fread(buffer, 1, 4, file);
            image[y][x].a = buffer[0];
            image[y][x].r = buffer[1];
            image[y][x].g = buffer[2];
            image[y][x].b = buffer[3];
        }
    }
\end{lstlisting}


Finally, we close the file and return the pointer to the matrix of pixels.


\section{Save .avs from memory}

The function \texttt{save\_avs} stores an image represented in memory as an AVS
file. Its declaration is:

\begin{lstlisting}
int save_avs(const char *filename, Pixel **image,
             int image_width, int image_height)
\end{lstlisting}

The function receives the output filename, a matrix of pixels, the image width,
and the image height. It returns \texttt{0} if the operation is successful and
\texttt{-1} if an error occurs.

\subsection{Opening the output file}

The file is opened in binary write mode:

\begin{lstlisting}
FILE *file = fopen(filename, "wb");
\end{lstlisting}

The \texttt{w} flag enables writing, while the \texttt{b} flag indicates that
the file must be treated as binary data. If the file cannot be created, the
function reports the error and returns \texttt{-1}.

\begin{lstlisting}
if (file == NULL) {
    printf("Could not create file.\n");
    return -1;
}
\end{lstlisting}

\subsection{Writing the image width}

The AVS format stores the width in the first four bytes of the file. These bytes
must be written in big-endian order. The bit-shift operators extract each byte
from the integer:

\begin{lstlisting}
buffer[0] = (unsigned char)((image_width >> 24) & 0xff);
buffer[1] = (unsigned char)((image_width >> 16) & 0xff);
buffer[2] = (unsigned char)((image_width >> 8) & 0xff);
buffer[3] = (unsigned char)(image_width & 0xff);
\end{lstlisting}

The four bytes are then written to the file using \texttt{fwrite}:

\begin{lstlisting}
if (fwrite(buffer, 1, 4, file) != 4) {
    fclose(file);
    return -1;
}
\end{lstlisting}

The return value of \texttt{fwrite} is checked to ensure that all four bytes
were written successfully.

\subsection{Writing the image height}

The height is written using the same big-endian conversion:

\begin{lstlisting}
buffer[0] = (unsigned char)((image_height >> 24) & 0xff);
buffer[1] = (unsigned char)((image_height >> 16) & 0xff);
buffer[2] = (unsigned char)((image_height >> 8) & 0xff);
buffer[3] = (unsigned char)(image_height & 0xff);
\end{lstlisting}

The resulting four bytes are written immediately after the width.
The first eight bytes of the AVS file contain the image dimensions:

\begin{verbatim}
4 bytes for the width
4 bytes for the height
\end{verbatim}

\subsection{Writing the pixel data}

After writing the header, the function iterates through every row and column of
the image:

\begin{lstlisting}
for (int y = 0; y < image_height; y++) {
    for (int x = 0; x < image_width; x++) {
\end{lstlisting}

Each pixel is stored in the order alpha, red, green, and blue:

\begin{lstlisting}
buffer[0] = image[y][x].a;
buffer[1] = image[y][x].r;
buffer[2] = image[y][x].g;
buffer[3] = image[y][x].b;
\end{lstlisting}

Each pixel occupies four bytes in the output file. The function writes those
bytes and verifies that the operation succeeds:

\begin{lstlisting}
if (fwrite(buffer, 1, 4, file) != 4) {
    fclose(file);
    return -1;
}
\end{lstlisting}

Pixels are written from left to right and from the top row to the bottom row.

\subsection{Closing the file}

Once the header and all pixel data have been written, the file is closed:

\begin{lstlisting}
if (fclose(file) != 0) {
    return -1;
}
\end{lstlisting}

If the file is closed successfully, the function returns zero:

\begin{lstlisting}
return 0;
\end{lstlisting}

\subsection{Using the function in \texttt{main}}

The image is first loaded into memory using \texttt{load\_avs}:

\begin{lstlisting}
Pixel **image = load_avs("image.avs");
\end{lstlisting}

The in-memory image is then saved as a new AVS file:

\begin{lstlisting}
if (save_avs("image_copy.avs", image, width, height) == 0) {
    printf("Image saved successfully.\n");
} else {
    printf("Could not save image.\n");
}
\end{lstlisting}

This creates \texttt{image\_copy.avs} with the same dimensions and pixel values
as the original image.

Finally, the memory allocated for the pixel matrix is released row by row:

\begin{lstlisting}
for (int y = 0; y < height; y++) {
    free(image[y]);
}
free(image);
\end{lstlisting}

This prevents memory leaks after the image has been saved.